

\section{La Surcharge Cognitive (Sweller) : Analyse + dictionnaire + traduction = blocage}
space{0.3cm}
oindent

L'enseignement du latin et du grec ancien traverse une crise épistémologique et pédagogique
profonde, souvent résumée par le constat d'une inefficacité structurelle à produire des lecteurs
fluides. Au cœur de cette problématique se trouve une méthodologie héritée du XIXe siècle, la
méthode Grammaire-Traduction (\textit{Grammar-Translation Method} ou GTM), qui repose sur une
séquence procédurale apparemment logique mais cognitivement coûteuse : face à un texte en langue
cible, l'apprenant est instruit d'analyser grammaticalement les formes (morphosyntaxe), de
rechercher le sens des mots inconnus dans un outil externe (dictionnaire), et de transposer le tout
dans sa langue maternelle (traduction).

L'hypothèse centrale de cette section est que l'équation suivante synthétise le blocage cognitif : \textit{analyser grammaticalement \+ chercher dans le dictionnaire \+ traduire \= surcharge de la mémoire de travail}. Cette équation ne relève pas d'une simple opinion pédagogique mais trouve une validation robuste dans les sciences cognitives modernes.

En mobilisant la Théorie de la Charge Cognitive (TCC) développée par John Sweller, ainsi que les modèles de la mémoire de travail d'Alan Baddeley, on peut démontrer que la juxtaposition de ces trois tâches sature les ressources limitées du cerveau humain, interdisant mécaniquement l'émergence de la compréhension sémantique profonde et de la fluidité de lecture.

L'analyse qui suit déconstruit les mécanismes neurocognitifs à l'œuvre. Elle explore comment
l'architecture de notre mémoire, conçue pour traiter l'information de manière sérielle et limitée,
s'effondre sous le poids de la charge intrinsèque des langues flexionnelles lorsqu'elle est aggravée
par la charge extrinsèque des procédures de la GTM. De l'effet d'attention partagée (\textit{Split-
Attention Effect}) induit par le dictionnaire à la compétition neuronale entre le traitement
métalinguistique et l'accès sémantique, nous établirons que l'accès au sens est physiologiquement
empêché par cette méthodologie. Enfin, nous examinerons les alternatives empiriquement validées,
telles que l'approche communicative de l'Institut Polis ou les méthodes de lecture extensive, qui
visent à restaurer l'homéostasie cognitive nécessaire à l'apprentissage.



\subsection{L'architecture cognitive humaine et ses contraintes : le cadre de sweller}

Pour comprendre la nature de l'échec induit par l'équation « Analyse \+ Dictionnaire \+ Traduction
», il est impératif de définir d'abord le terrain sur lequel se joue l'apprentissage :
l'architecture cognitive humaine. Les travaux de John Sweller, initiés dans les années 1980 et
affinés jusqu'à aujourd'hui, fournissent le cadre théorique indispensable : la Théorie de la Charge
Cognitive (TCC).



\subsection{La mémoire de travail : le goulot d'étranglement du traitement de l'information}

Au centre de la cognition se trouve la mémoire de travail (MT). Contrairement à la mémoire à long
terme (MLT), qui possède une capacité de stockage théoriquement illimitée, la MT est sévèrement
contrainte. Elle est le lieu du traitement conscient, là où l'information est manipulée, organisée
et comprise.\cite{II-1-3-ref1}

Selon le modèle multicompartimental de Baddeley, souvent intégré aux analyses de la charge
cognitive, la MT se compose de plusieurs sous-systèmes :

1. \textbf{L'Administrateur Central (\textit{Central Executive}) :} Il pilote l'attention, coordonne
les tâches et gère les conflits cognitifs. C'est lui qui est sollicité lors de l'analyse
grammaticale consciente et de la prise de décision en traduction.\cite{II-1-3-ref2}

2. \textbf{La Boucle Phonologique (\textit{Phonological Loop}) :} Elle maintient temporairement les
informations verbales et auditives par un processus de répétition mentale. Dans la lecture, elle est
cruciale pour garder en mémoire le début d'une phrase le temps d'arriver à la fin.\cite{II-1-3-ref2}

3. \textbf{Le Calepin Visuo-spatial (\textit{Visuospatial Sketchpad}) :} Il traite les informations
visuelles et spatiales.

4. \textbf{Le Buffer Épisodique :} Il assure l'interface avec la mémoire à long terme.

La capacité de la MT est limitée à environ 7 (±2) éléments d'information, voire 4 pour des éléments
complexes, et sa durée de rétention sans répétition est de quelques secondes.\cite{II-1-3-ref1}
Toute tâche d'apprentissage qui exige de maintenir et de manipuler simultanément plus d'éléments que
cette capacité ne le permet entraîne une surcharge cognitive (\textit{cognitive overload}). Lorsque
la surcharge survient, le traitement de l'information s'arrête, l'apprentissage échoue et la
compréhension est bloquée.\cite{II-1-3-ref3}



\subsection{La typologie de la charge cognitive selon sweller}

Sweller distingue trois types de charge qui s'additionnent pour former la charge cognitive totale.
Cette distinction est fondamentale pour diagnostiquer les défauts de la méthode
GTM.\cite{II-1-3-ref3}

\begin{table}[h!]\small\centering\begin{tabular}{| p{2cm} | p{2.2cm} | p{2.2cm} | p{2.2cm} |}\hline\textbf{Type de Charge} & \textbf{Définition Cognitive} & \textbf{Manifestation dans l'Apprentissage du Latin/Grec} & \textbf{Impact de la Méthode GTM} \\ \hline\textbf{Charge Intrinsèque} (\textit{Intrinsic Load}) & Liée à la complexité inhérente du matériel et à l'interactivité des éléments (\textit{element interactivity}). Elle dépend de l'expertise de l'apprenant. & La morphologie flexionnelle (déclinaisons, conjugaisons), l'ordre des mots libre (hyperbates), la complexité syntaxique. & \textbf{Élevée et incompressible.} La nature même des langues anciennes impose une charge intrinsèque forte pour le novice.\cite{II-1-3-ref7} \\ \hline\textbf{Charge Extrinsèque} (\textit{Extraneous Load}) & Générée par la manière dont l'information est présentée ou par les activités imposées à l'apprenant. Elle ne contribue pas à l'apprentissage (construction de schémas). & La nécessité de manipuler un dictionnaire, de consulter des tableaux de grammaire séparés, de traduire mentalement en L1. & \textbf{Massive et artificielle.} La GTM ajoute délibérément des obstacles procéduraux qui ne sont pas la langue elle-même.\cite{II-1-3-ref5} \\ \hline\textbf{Charge Germane} (\textit{Germane Load}) & Ressources de la MT consacrées à l'apprentissage effectif : la construction et l'automatisation de schémas en mémoire à long terme. & L'effort pour comprendre le sens, relier le mot \textit{pater} au concept de \og père\fg{}, intégrer la syntaxe. & \textbf{Réduite à néant.} Si (Intrinsèque \+ Extrinsèque) \> Capacité totale, il ne reste aucune ressource pour la charge germane.\cite{II-1-3-ref3} \\ \hline\end{tabular}\end{table}L'équation critique de notre recherche peut donc se réécrire en termes swelleriens :

$$Charge Total \= (Analyse Grammaticale \[Intrinsèque mal gérée\]) \+ (Dictionnaire \[Extrinsèque\])
\+ (Traduction \[Extrinsèque\])$$

Si ce total dépasse la capacité de la mémoire de travail, l'accès au sens est, par définition,
impossible.



\subsection{L'interactivité des éléments : pourquoi le latin est "lourd"}

Le concept d'« interactivité des éléments » est central chez Sweller. Une tâche a une faible
interactivité si les éléments peuvent être appris isolément (ex: apprendre une liste de vocabulaire
sans contexte). Elle a une haute interactivité si les éléments dépendent les uns des autres pour
être compris.\cite{II-1-3-ref3}

Dans une phrase latine comme \textit{« Magnam diserti poetae filiam amat »}, l'apprenant ne peut pas
traiter \textit{Magnam} (accusatif féminin singulier) isolément. Il doit le maintenir en mémoire de
travail jusqu'à trouver \textit{filiam} (accusatif féminin singulier), tout en traitant
\textit{diserti poetae} (génitif masculin singulier) intercalé, et \textit{amat} (verbe final).
Cette dépendance mutuelle des mots crée une charge intrinsèque explosive. Si la méthode
d'enseignement ajoute des charges extrinsèques (chercher \textit{diserti} dans le dictionnaire,
analyser explicitement la terminaison \textit{\-ae}), le système sature
instantanément.\cite{II-1-3-ref7}



\subsection{Analyse de la composante 1 : l'analyse grammaticale comme tâche concurrente et perturbatrice}

La première injonction de la méthode traditionnelle est d'« analyser » : identifier consciemment les
cas, les temps et les modes avant de tenter de construire du sens. Cette approche, souvent défendue
comme rigoureuse, constitue en réalité une aberration cognitive pour le novice.



\subsection{Le coût cognitif du « parsing » conscient}

Le traitement syntaxique, ou \textit{parsing}, peut s'effectuer de deux manières : implicite
(automatique) ou explicite (consciente). Chez le locuteur natif ou l'expert, le parsing est un
processus rapide, encapsulé et peu coûteux en ressources attentionnelles. Le cerveau détecte les
structures grammaticales sans effort conscient.\cite{II-1-3-ref8}

La méthode GTM, en revanche, impose un \textbf{parsing explicite}. L'élève doit observer la
terminaison, récupérer consciemment une règle déclarative (ex: \og la terminaison \-is peut être
datif ou ablatif pluriel\fg{}), et tester des hypothèses.

\begin{itemize}\item \textbf{Saturation de l'Administrateur Central :} Les études de Deida Perea (2020) et Juffs (2005) montrent que le parsing grammatical conscient mobilise massivement l'administrateur central de la mémoire de travail.\cite{II-1-3-ref10} L'attention est détournée du \textit{message} (le sens) vers le \textit{code} (la forme).\item \textbf{Compétition des Ressources :} La mémoire de travail étant limitée, chaque unité d'attention allouée à l'analyse morphologique est une unité soustraite à la construction du modèle sémantique. Comme le souligne Bextermöller (2018), les stratégies de bas niveau (décodage) monopolisent les ressources, empêchant les processus de haut niveau (intégration du sens, inférence).\cite{II-1-3-ref11}\end{itemize}

\subsection{L'hypothèse du moniteur de krashen et la lenteur de traitement}

Stephen Krashen a théorisé ce phénomène avec l'Hypothèse du Moniteur (Monitor Hypothesis). Selon
Krashen, l'apprentissage conscient des règles (Learning) ne sert qu'à \og surveiller\fg{} ou
corriger la production, mais ne génère pas la fluidité.\cite{II-1-3-ref13}

L'utilisation du \og Moniteur\fg{} (l'analyse grammaticale consciente) requiert trois conditions :
du temps, une focalisation sur la forme, et la connaissance de la règle.

\begin{itemize}\item \textbf{Le Facteur Temps :} L'analyse grammaticale est intrinsèquement lente. Elle oblige à une lecture mot à mot, brisant le flux nécessaire à la compréhension globale.\item \textbf{L'Effondrement de la Compréhension :} En ralentissant le traitement à une vitesse inférieure au seuil de dégradation de la boucle phonologique (environ 2 secondes sans répétition), l'analyse grammaticale provoque l'oubli du début de la phrase avant que la fin ne soit atteinte.\cite{II-1-3-ref2} L'élève \og analyse\fg{} correctement chaque mot mais ne comprend pas la phrase, car les morceaux ne peuvent plus être assemblés en un tout cohérent dans le \textit{buffer} épisodique.\end{itemize}

\subsection{Neurobiologie du traitement grammatical}

Les recherches en neurosciences confirment que le traitement grammatical automatisé et l'analyse
consciente sollicitent des réseaux neuronaux distincts. Le traitement procédural (fluide) repose
notamment sur les ganglions de la base et l'aire de Broca (BA 44) pour le traitement
hiérarchique.\cite{II-1-3-ref8}

En revanche, l'analyse grammaticale de type « résolution de problème » (telle qu'imposée par la GTM)
active des réseaux fronto-pariétaux associés à l'effort cognitif intense et au contrôle
attentionnel.\cite{II-1-3-ref8} Cette activation corticale massive est énergivore et fatigue
rapidement l'apprenant, menant à une baisse de performance (épuisement cognitif) bien plus rapide
que lors d'une lecture fluide.



\subsection{L'erreur de l'enseignement explicite isolé}

L'enseignement explicite de la grammaire, s'il n'est pas immédiatement suivi d'une pratique massive
permettant la \og proceduralisation\fg{} (Anderson), reste une connaissance déclarative inerte. Les
études montrent que la connaissance explicite des règles ne garantit pas leur application correcte
en temps réel.\cite{II-1-3-ref16} Pire, l'effort pour récupérer ces règles en mémoire à long terme
(MLT) pendant la lecture ajoute une charge cognitive supplémentaire. Au lieu de \textit{reconnaître}
une forme (processus rapide), l'élève doit la \textit{reconstruire} (processus lent et
coûteux).\cite{II-1-3-ref17}



\subsection{Analyse de la composante 2 : le dictionnaire et la violence de l'effet d'attention partagée}

La deuxième variable de l'équation, « chercher dans le dico », est identifiée dans la littérature
psychopédagogique comme une source majeure de charge extrinsèque, connue sous le nom d'Effet
d'Attention Partagée (\textit{Split-Attention Effect}).



\subsection{Le mécanisme de l'attention partagée}

L'effet d'attention partagée survient lorsque l'apprenant doit intégrer mentalement des informations
provenant de deux sources physiquement ou temporellement séparées pour comprendre un
concept.\cite{II-1-3-ref19} Dans le cas de la lecture avec dictionnaire :

1. \textbf{Source A :} Le mot inconnu dans le texte latin/grec.

2. \textbf{Source B :} La définition dans le dictionnaire (livre papier, fin du manuel, ou outil
numérique).

Le processus cognitif imposé est dévastateur pour la mémoire de travail :

\begin{itemize}\item \textbf{Rupture Visuelle :} L'œil quitte le texte.\item \textbf{Maintien Coûteux :} L'apprenant doit garder l'orthographe du mot en mémoire de travail tout en effectuant une tâche de recherche visuelle (scan alphabétique) dans le dictionnaire.\item \textbf{Interférence :} La recherche visuelle dans le dictionnaire introduit de nouvelles informations (autres mots, définitions multiples) qui entrent en compétition avec les données maintenues en mémoire.\cite{II-1-3-ref20}\item \textbf{Réintégration Pénible :} Une fois la définition trouvée, l'apprenant doit revenir au texte, retrouver sa ligne, et tenter d'insérer le sens trouvé dans le contexte. Souvent, le contexte a été perdu (oublié) à cause du délai et de l'interférence, obligeant à une relecture.\cite{II-1-3-ref21}\end{itemize}

\subsection{Preuves empiriques de l'effet négatif du dictionnaire}

Les travaux de Yeung et al. (1997) ont démontré que les apprenants lisant avec des définitions
intégrées (glosses en marge ou interlinéaires) obtiennent de meilleurs résultats en compréhension
que ceux utilisant une liste de vocabulaire séparée.\cite{II-1-3-ref22}

De même, Al-Shehri \& Gitsaki (2010) ont montré que même avec des dictionnaires en ligne (plus
rapides), l'effet de rupture subsiste. La charge cognitive générée par la gestion de l'outil externe
consomme des ressources qui ne sont plus disponibles pour la compréhension
profonde.\cite{II-1-3-ref23}

L'étude de Mayer et Moreno sur l'apprentissage multimédia confirme que la proximité spatiale
(Contiguity Principle) est essentielle. Séparer le mot de son sens (texte vs dictionnaire) viole ce
principe et impose une charge extrinsèque inutile.\cite{II-1-3-ref19}



\subsection{Le seuil lexical de 95-98\% : pourquoi le dictionnaire est une impasse}

Les recherches de Paul Nation et Batia Laufer sur les seuils lexicaux apportent un éclairage
quantitatif crucial. Pour qu'une lecture soit fonctionnelle et permette l'accès au sens sans
surcharge excessive, le lecteur doit connaître \textbf{95\% à 98\%} des mots du
texte.\cite{II-1-3-ref25}

\begin{itemize}\item \textbf{À 98\% (1 mot inconnu sur 50) :} La lecture est fluide, le sens global est clair, et le mot inconnu peut être deviné par inférence contextuelle.\item \textbf{À 95\% (1 mot inconnu sur 20) :} La compréhension est possible mais demande un effort.\item \textbf{En dessous de 95\% :} L'accès au sens s'effondre. Le recours au dictionnaire devient si fréquent (tous les 2 ou 3 mots) que la tâche change de nature : ce n'est plus de la lecture, c'est du déchiffrage (\textit{decoding}).\end{itemize}Dans l'enseignement traditionnel, on demande souvent aux élèves de traduire des textes (César,
Cicéron) où ils ne maîtrisent parfois que 60\% ou 70\% du vocabulaire. Ils sont donc bien en deçà du
\og seuil de rupture\fg{}. Le dictionnaire devient une béquille permanente qui empêche la marche. La
surcharge est constante, et l'accès au sens global est impossible car la mémoire de travail est
saturée par la gestion des lacunes lexicales.\cite{II-1-3-ref28}



\subsection{Analyse de la composante 3 : « traduire » ou le coût du transcodage mental}

La troisième composante, la traduction systématique (version), transforme l'acte de lecture en un
exercice de résolution de problèmes complexes, ajoutant une couche finale de charge cognitive.



\subsection{Distinction entre lecture et traduction}

La psycholinguistique distingue fermement la \textit{lecture} (\textit{reading for comprehension})
de la \textit{traduction} (\textit{translation}).

\begin{itemize}\item \textbf{Lecture (Information Processing) :} C'est un processus vertical où les formes linguistiques (mots, phrases) activent des représentations sémantiques (concepts, images mentales). L'objectif est le sens. C'est un processus qui gagne à être automatisé.\cite{II-1-3-ref12}\item \textbf{Traduction (Problem Solving) :} C'est un processus horizontal et itératif. Il faut non seulement comprendre le sens (déverbalisation), mais ensuite le ré-encoder dans une autre langue (L1) avec ses propres contraintes syntaxiques et lexicales. C'est une tâche de \og résolution de problème\fg{} à haute exigence cognitive.\cite{II-1-3-ref31}\end{itemize}

\subsection{Le fardeau cognitif du transcodage}

La méthode GTM exige la traduction non pas comme \textit{résultat} de la compréhension, mais comme
\textit{moyen} d'y accéder. C'est une erreur pédagogique majeure.

1. \textbf{Double Gestion des Systèmes :} L'élève doit maintenir actifs deux systèmes linguistiques
en parallèle (L2 et L1). Cela exige un contrôle inhibiteur constant pour éviter les interférences
(faux amis, calques syntaxiques), ce qui épuise l'administrateur central.\cite{II-1-3-ref33}

2. \textbf{Surcharge de la Mémoire de Travail :} Au lieu de mapper le mot latin \textit{canis}
directement vers le concept {CHIEN}, l'élève mappe \textit{canis} vers le mot français « chien ». Ce
pas supplémentaire, multiplié par chaque mot de la phrase, augmente la latence et la charge.

3. \textbf{Perte du Sens Global :} Focalisé sur la recherche de l'équivalent français précis mot à
mot (transcodage), l'élève perd souvent de vue la macro-structure du texte. Il produit une suite de
mots français grammaticalement correcte mais dont le sens global lui échappe parfois totalement («
J'ai traduit mais je ne sais pas ce que ça veut dire »).\cite{II-1-3-ref4}



\subsection{L'inefficacité de la traduction pour l'acquisition}

Les recherches en Acquisition des Langues Secondes (SLA) indiquent que la traduction, bien qu'utile
comme exercice ponctuel de vérification, est un moyen inefficace d'acquisition. Elle ne favorise pas
le développement de l'interlangue ni l'automatisation des processus de lecture en L2. Elle maintient
l'apprenant dans une dépendance à la L1, empêchant le développement d'une compétence de lecture
autonome.\cite{II-1-3-ref36} L'étude de Humphreys (2020) suggère que sans une capacité de penser
directement dans la langue cible, la fluidité de lecture réelle est
inatteignable.\cite{II-1-3-ref38}



\subsection{Synthèse : la tempête parfaite de la surcharge cognitive}

L'équation posée en II.1.\cite{II-1-3-ref3} se révèle être une description précise d'un effondrement
cognitif programmé. En additionnant les charges, nous obtenons un dépassement systématique de la
capacité de traitement.



\subsection{Modélisation de la surcharge}

Si nous représentons la capacité de la Mémoire de Travail (MT) par un indice fictif de 100 unités :

\begin{itemize}\item \textbf{Charge Intrinsèque (Latin/Grec) :} \~40 unités. La morphologie complexe, l'ordre des mots, l'absence d'articles exigent déjà une attention soutenue.\item \textbf{Charge Extrinsèque 1 (Analyse grammaticale) :} \~30 unités. Le rappel conscient des règles et des tableaux de déclinaison.\item \textbf{Charge Extrinsèque 2 (Dictionnaire \- Split Attention) :} \~40 unités. La rupture visuelle, le maintien en boucle, la recherche, la réintégration.\item \textbf{Charge Extrinsèque 3 (Traduction \- Transcodage) :} \~30 unités. La gestion de la L1 et du mapping lexical.\end{itemize}\textbf{Conséquence :} La Charge Germane (ressources pour comprendre et apprendre) est négative ou
nulle. L'apprentissage ne peut pas avoir lieu. L'apprenant est en mode \og survie cognitive\fg{},
traitant des fragments isolés sans jamais pouvoir construire la cohérence
textuelle.\cite{II-1-3-ref3}



\subsection{L'empêchement de l'accès au sens}

L'accès au sens nécessite la construction d'un \og modèle de situation\fg{} (Kintsch). Ce modèle se
construit par l'intégration progressive des propositions sémantiques.

\begin{itemize}\item La \textbf{surcharge de la MT} efface les propositions précédentes avant qu'elles ne soient intégrées.\item L'\textbf{analyse grammaticale} morcelle le texte en unités trop petites (morphèmes) pour porter du sens.\item Le \textbf{dictionnaire} interrompt la continuité temporelle nécessaire à l'intégration.\item La \textbf{traduction} substitue un exercice de code à un exercice de sens.\end{itemize}Le résultat est une \og cécité sémantique\fg{}. L'élève peut identifier un ablatif absolu, trouver
le sens des mots dans le dictionnaire, et écrire une traduction française, tout en n'ayant aucune
représentation mentale de la scène décrite par le texte. L'accès au sens est structurellement
empêché.\cite{II-1-3-ref40}



\subsection{Vers une résolution de l'équation : alternatives pédagogiques fondées sur la preuve}

Pour restaurer l'accès au sens, il est nécessaire de réduire drastiquement la charge cognitive
extrinsèque et de gérer la charge intrinsèque. La littérature scientifique et les pratiques
innovantes offrent des solutions validées.



\subsection{Éliminer l'effet d'attention partagée : textes glossés et intégrés}

La solution la plus immédiate à la surcharge du dictionnaire est l'utilisation de textes avec
vocabulaire intégré (en marge, en bas de page, ou interlinéaire).

\begin{itemize}\item \textbf{Preuve :} Yeung (1997) et d'autres ont montré que cela libère des ressources cognitives pour la compréhension.\cite{II-1-3-ref22}\item \textbf{Application :} Les éditions de type \og Tiered Readers\fg{} ou les textes annotés permettent de maintenir le flux visuel et attentionnel sur le texte.\cite{II-1-3-ref42}\end{itemize}

\subsection{Réduire la charge grammaticale : l'approche inductive et l'input compréhensible}

Au lieu de l'analyse explicite \textit{a priori}, les méthodes basées sur l'Input Compréhensible
(CI) favorisent une acquisition implicite.

\begin{itemize}\item \textbf{Méthode Naturelle (Ørberg) :} Le manuel \textit{Lingua Latina Per Se Illustrata} introduit la grammaire par le contexte et la répétition, permettant une induction des règles sans surcharge métalinguistique. L'élève \textit{voit} le fonctionnement de l'accusatif avant de le nommer.\cite{II-1-3-ref44}\item \textbf{Preuve :} L'apprentissage implicite permet de traiter la grammaire de manière procédurale (automatique), réduisant la charge sur l'administrateur central.\cite{II-1-3-ref18}\end{itemize}

\subsection{Remplacer la traduction par la compréhension directe : les méthodes actives}

Les approches communicatives ou \og vivantes\fg{} (Living Latin/Greek), défendues par l'Institut
Polis (Christophe Rico) ou des pédagogues comme Patrick Owens et Justin Slocum Bailey, visent à
court-circuiter la L1.

\begin{itemize}\item \textbf{Méthode Polis :} L'immersion totale et l'usage de l'oral (koineisation) forcent le cerveau à traiter la langue cible directement. L'oralisation aide à internaliser les structures (boucle phonologique efficace) et à développer une fluidité qui soutient la lecture.\cite{II-1-3-ref38}\item \textbf{Total Physical Response (TPR) :} Associer le mot à une action physique ancre le sens sans passer par la traduction, utilisant la mémoire procédurale et réduisant la charge de la MT.\cite{II-1-3-ref49}\item \textbf{Fluidité de Lecture :} L'objectif est d'atteindre une vitesse de lecture suffisante (au moins 100-150 mots/minute) pour que la mémoire de travail puisse maintenir la cohérence du discours. Cela nécessite une lecture extensive de textes accessibles (niveau i+1 de Krashen) plutôt que le décryptage intensif de textes trop difficiles.\cite{II-1-3-ref51}\end{itemize}

\subsection{Tableau comparatif des charges cognitives}

\begin{table}[h!]\small\centering\begin{tabular}{| p{2cm} | p{2.2cm} | p{2.2cm} | p{2.2cm} |}\hline\textbf{Composante} & \textbf{Méthode Grammaire-Traduction (GTM)} & \textbf{Approche Communicative / Input Compréhensible} & \textbf{Gain Cognitif} \\ \hline\textbf{Grammaire} & Analyse explicite, déductive. Charge : \textbf{LOURDE} & Acquisition implicite, inductive, contextuelle. Charge : \textbf{MODÉRÉE} & Libération de l'Administrateur Central. \\ \hline\textbf{Vocabulaire} & Dictionnaire externe (Split-Attention). Charge : \textbf{TRÈS LOURDE} & Glosses intégrées, vocabulaire fréquentiel, gestes. Charge : \textbf{FAIBLE} & Maintien de la Boucle Phonologique. \\ \hline\textbf{Sémantique} & Traduction (Transcodage L2 \-\> L1). Charge : \textbf{LOURDE} & Compréhension directe (L2 \-\> Sens). Charge : \textbf{NATURELLE} & Élimination de l'inhibition L1. \\ \hline\textbf{Résultat} & \textbf{SURCHARGE} (Échec du sens) & \textbf{FLUIDITÉ} (Accès au sens) & Apprentissage durable (Germane). \\ \hline\end{tabular}\end{table}

\subsection{Conclusion}

L'équation \textit{« Analyse Grammaticale \+ Recherche Dictionnaire \+ Traduction \= Surcharge
Cognitive »} est une réalité neurocognitive incontestable. Elle décrit avec précision pourquoi la
méthode traditionnelle échoue à produire des lecteurs compétents en langues anciennes, malgré des
années d'efforts. En forçant l'apprenant à effectuer simultanément trois tâches à haute charge
cognitive qui entrent en compétition pour les mêmes ressources neuronales limitées, cette
méthodologie sature la mémoire de travail et bloque physiologiquement les processus d'intégration
sémantique.

L'accès au sens n'est pas une récompense que l'on obtient \textit{après} l'analyse ; c'est un
processus qui doit être facilité \textit{pendant} la lecture. Les données de la recherche, des
travaux de Sweller sur la charge cognitive aux études oculométriques de Bextermöller sur la lecture
du latin, convergent vers une conclusion unique : pour enseigner efficacement les langues anciennes,
il faut alléger la charge extrinsèque. Cela implique d'abandonner le dictionnaire externe au profit
de textes glossés, de remplacer l'analyse explicite prématurée par une exposition inductive massive,
et de substituer la traduction systématique par des pratiques de compréhension directe et active. Ce
n'est qu'en respectant les contraintes biologiques de notre architecture cognitive que l'on pourra
transformer le déchiffrage laborieux en lecture véritable.

